\section{Les 6 -- Irreduceerbare elementen, restklassenringen en eindige velden}

\begin{roundbox}[CornflowerBlue!10][RoyalPurple]
    \textbf{Bézout + Euclidische deling geven je algoritmes voor inverses, ggd's én voor het bouwen van velden.}

    \vspace{0.25cm}

    Eénzelfde redening ($1 = sa + tb$) zal deze les vier keer opduiken: bij inverses modulo $n$, bij het priemlemma, bij eindige velden $\mathbb{F}_p$, en bij velduitbreidingen $F[x]/(g)$.
\end{roundbox}

\subsection*{Recap van vorige les}

Een \inlinetext{Euclidische ring} is een integriteitsdomein $D$ met een "\emph{grootte}"-functie
\begin{equation*}
    d : D_0 \to \mathbb{N}
\end{equation*}
zodat:
\begin{enumerate}
    \item $\forall a,b \in D\backslash \{0\}$ geldt $d(a) \le d(ab)$, en
    \item $\forall a \in D$ en $\forall b \in D \backslash \{0\}$ geldt \textbf{deling met rest}. Maw. bestaan $q,\ r \in D$ zodat $a = qb + r$ met $r=0$ of $d(r) < d(b)$
\end{enumerate}
Omdat $d(r)$ strikt daalt en in $\mathbb{N}$ leeft, moet herhaald delen uiteindelijk stoppen -- daarom ook dat het \textbf{algoritme van Euclides} werkt, en dat \textbf{Bézout} garandeert dat $\operatorname{ggd}(a,\ b) = sa + tb$.

\subsubsection*{Algoritme van Euclides}

\begin{nicealgorithm}*[Algoritme van Euclides]
    \nextpart[]*
    Voor twee elementen $a,\ b$ in een Euclidische ring berekent het algoritme van Euclides de grootste gemene deler door herhaald deling met rest uit te voeren:
    \begin{equation*}
        \begin{aligned}
            a   &= q_1b + r_1 \\
            b   &= q_2r_1 + r_2 \\
            r_1 &= q_3r_2 + r_3 \\
                &\;\; \vdots \\
            r_{k-2} &= q_kr_{k-1} + r_k \\
            r_{k-1} &= q_{k+1}r_k + 0
        \end{aligned}
    \end{equation*}
    De laatste niet-nul rest $r_k$ is dan
    \begin{equation*}
        r_k = \operatorname{ggd}(a,\ b)
    \end{equation*}
\end{nicealgorithm}

\subsubsection*{Stelling van Bézout}

In een Euclidische ring kan de ggd van $a$ en $b$ beschreven worden als een lineaire combinatie van $a$ en $b$:
\begin{equation*}
    \operatorname{ggd}(a,\ b) = sa + tb
\end{equation*}
voor zekere $s,\ t \in D$.

Als $a$ en $b$ onderling ondeelbaar zijn, dus $\operatorname{ggd}(a,\ b) = 1$, dan bestaan $s,\ t$ zodat
\begin{equation*}
    1 = sa + tb
\end{equation*}
Dit is de vorm die later heel vaak gebruikt wordt om inverses te vinden.

\subsection{Veeltermen over een veld: \texorpdfstring{$F[x]$}{F[x]} als Euclidische ring}

Dit is de grote uitbreiding van deze les: de Euclidische theorie geldt niet enkel voor $\mathbb{Z}$, maar ook voor \textbf{veeltermen}.

Als $F$ een veld is, dan is de veeltermring $F[x]$ een Euclidische ring met als grootte-functie de graad:
\begin{equation*}
    d(f) = \operatorname{def}(f)
\end{equation*}
Hier is $F[x]$ de ring van univariate veetermen:
\begin{equation*}
    F[x] = \{a_0 + a_1x + \dots + a_nx^n \mid a_i \in F\}
\end{equation*}
De Euclidische deling is dan gewoon staartdeling van veeltermen.

\subsubsection{Waarom moet \texorpdfstring{$F$}{F} een veld zijn?}

Bij het delen wil je de hoogste term van $f$ wegwerken met een veelvoud van $g$. Stel:
\begin{equation*}
    f(x) = a_nx^n + \dots,\ \qquad g(x) = b_mx^m + \dots
\end{equation*}
waar $a_n,\ b_m \neq 0$ en $ n \ge m$. Dan trek je af:
\begin{equation*}
    \frac{a_n}{b_m}x^{n-m}g(x)
\end{equation*}
Dat vereist dat je door $b_m$ kunt \textbf{delen}, dus dat $b_m$ inverteerbaar is. Daarom moeten de coëfficiënten dus in een \textbf{veld} liggen!

\textbf{\underbar{Voorbeeld:}} ggd van veeltermen in $\mathbb{Q}[x]$.\newline
Bepaal:
\begin{equation*}
    \operatorname{ggd}(x^3 + x^2 - 2x,\ x^2 - 1)
\end{equation*}
\underbar{Oplossing:} Neem
\begin{equation*}
    f = x^3 + x^2 - 2x, \qquad g = x^2 - 1 
\end{equation*}
Deel $f$ door $g$:
\begin{equation*}
    (x+1)(x^2-1) = x^3 + x^2 - x - 1
\end{equation*}
Dus:
\begin{equation*}
    f - (x+1)g = (x^3 + x^2 - 2x) - (x^3 + x^2 - x - 1) = 1 - x
\end{equation*}
De rest is dus $1-x$, wat geassocieerd is met $x-1$. Volgende stap:
\begin{equation*}
    x^2 - 1 = (x+1)(x-1) + 0
\end{equation*}
Dus de ggd is, op eenheid na, $x-1$. Omdat men bij veeltermen meestal de \textbf{monische} ggd kiest:
\begin{equation*}
    \operatorname{ggd}(x^3 + x^2 - 2x,\ x^2 - 1)
\end{equation*}

\begin{barbox}[WildStrawberry][gray!3]
    \textcolor{WildStrawberry}{\textbf{Let op:}} In $F[x]$ zijn alle niet-nul constante veeltermen eenheden ($F[x]^{\times} = F \backslash \{0\}$). Daardoor is de ggd maar bepaald op een constante factor na. Afspraak: kies de \textcolor{WildStrawberry}{\textbf{monische}} ggd (leidende coëfficiënt $1$).
\end{barbox}

\textbf{\underbar{Herinner:}} Een element $u$ in een ring $A$ heet een \inlinetext{eenheid} als het inverteerbaar is voor de vermenigvuldiging. De verzameling van eenheden noteren we dan als:
\begin{equation*}
    A^{\times} = \{a \in A \mid a \text{ is inverteerbaar}\}
\end{equation*}

\subsection{Het priemlemma}

Dit lemma is de sleutel tot unieke factorisatie én tot de veldconstructies later.

We hebben dit al eens gedefinieerd (\Cref{H4_irreduceerbaar_element}), maar we doen het nog eens ter opfrissing.

\begin{nicedefinition}*[Irreduceerbaar element = Priemelement]
    \nextpart[]*
    Zij $\langle A, +, \cdot \rangle$ een integriteitsdomein (dus een commutatieve ring zonder nuldelers) en zij $p \in A\backslash\{0\}$. Dan noemen we $p$ een \textcolor{PineGreen!75!black}{\textbf{irreduceerbaar = priemelement}} als:
    \begin{enumerate}
        \item $p$ is geen eenheid (met eenheid bedoelen we elementen die inverteerbaar zijn). Dus
        \begin{equation*}
            p \notin A^{\times} = \{a \in A \mid \exists b \in A : ab=1\}
        \end{equation*} 
        \item Als $p=ab$ met $a,b \in A$ , dan is $a$ of $b$ een eenheid ($a \in A^{\times}$ of $b \in A^{\times}$)
    \end{enumerate}
\end{nicedefinition}

\textbf{\underbar{Merk op:}} Dit moet gelden voor elke ontbinding $ab$ van $p$!

\newpage
\begin{nicelemma}*[Priemlemma]
    \nextpart[]*
    In een Euclidische ring geld: Als $p$ irreduceerbaar is en
    \begin{equation*}
        p \mid ab,
    \end{equation*}
    dan geldt
    \begin{equation*}
        p\mid a \quad \text{ of } \quad p \mid b
    \end{equation*}
\end{nicelemma}

\inlinetext{\underbar{Bewijs:}}[Goldenrod!44] Stel dat $p \mid ab$. We moeten dan aantonen dat:
\begin{equation*}
    p \mid a \qquad \text{of} \qquad p \mid b
\end{equation*}
Als $p \mid a$, dan zijn we klaar. Veronderstel nu dat
\begin{equation*}
    p \nmid a
\end{equation*}
Dan moet je tonen dat het noodzakelijk is dat $p \mid b$. Beschouw de ggd:
\begin{equation*}
    \operatorname{ggd}(p,\ a)
\end{equation*}
Omdat $p$ irreduceerbaar is, zijn de enige delers van $p$ eenheden of elementen die geassocieerd zijn met $p$. Aangezien $\operatorname{ggd}(p,\ a)$ een deler is van $p$, zijn er dus twee mogelijkheden. Neem $d = \operatorname{ggd}(p,\ a)$. Dan geldt per definitie: $d \mid p$ en $d \mid a$. Omdat $d \mid p$, bestaat er een $c \in A$ zodat:
\begin{equation*}
    p = dc
\end{equation*}
Maar omdat $p$ irreduceerbaar is, is ofwel $d$ een eenheid, ofwel $c$ een eenheid. Als $d$ een eenheid is zijn we klaar. Neem nu dat $c$ een eenheid is, dan geld:
\begin{equation*}
    \begin{aligned}
        p &= dc \\
        &\Downarrow \textcolor{Green}{\text{ c inverteerbaar}} \\
        d &= pc^{-1} \\
    \end{aligned}
\end{equation*}
Maar $d\mid a$, dus er bestaat een $k \in A$ zodat:
\begin{equation*}
    \begin{aligned}
        a &= dk \\
          &= (pc^{-1})k = p(c^{-1}k)
    \end{aligned}
\end{equation*}
Dus $p \mid a$, maar we hadden aangenomen dat $p \nmid a$.
\begin{equation*}
    \begin{aligned}
        &\operatorname{ggd}(p,\ a) \text{ is een eenheid} \\
        &\operatorname{ggd}(p,\ a) \text{ is geassocieerd met } p
    \end{aligned}
\end{equation*}
Maar als $\operatorname{ggd}(p,a)$ geassocieerd zou zijn met $p$, dan zou $p$ ook $a$ delen. Dat is in tegenspraak met onze aanname $p \nmid a$. Dus moet gelden:
\begin{equation*}
    \operatorname{ggd}(p,\ a) = 1
\end{equation*}
tot op vermenigvuldiging met een eenheid. Met andere woorden: $p$ en $a$ zijn relatief priem. Omdat we in een Euclidische ring werken, geldt de stelling van Bézout. Er bestaan dus $s,t\in D$ zodat
\begin{equation*}
    \begin{aligned}
        1 &= sp + ta \\
          &\Downarrow \textcolor{Green}{\;\; \cdot b} \\
        b &= (sp + ta) \\
          &\Downarrow \textcolor{Green}{\text{ distributiviteit}} \\
        b &= spb + tab  
    \end{aligned}
\end{equation*}
Nu geldt: $p \mid spb$ (want $spb = p(sb)$). Ook geldt: $p \mid tab$ want per aanname is $p \mid ab$.
We hebben dus dat $p$ beide termen van de rechterkant deelt:
\begin{equation*}
    p \mid spd \qquad \text{en} \qquad p \mid tab
\end{equation*}
Daarom deelt $p$ ook hun som:
\begin{equation*}
    p \mid (spb + tab)
\end{equation*}
Maar $spb + tab = b$, dus $p \mid b$! We hebben nu dus aangetoond dat
\begin{equation*}
    p \nmid a \implies p \mid b
\end{equation*}
en dus is het lemma bewezen.

\subsubsection{Unieke factorisatie}

In een Euclidische ring heeft elk niet-nul element een \inlinetext{factorisatie in irreduceerbare elementen}[HighlightColor!66], uniek op volgorde en eenheden na.
\begin{equation*}
    a = u\;p_1\;p_2\;\dots\;p_r
\end{equation*}
waar $u$ een eenheid is en de $p_i$ irreduceerbaar zijn. Als ook geldt dat:
\begin{equation*}
    a = v\;q_1\;q_2\;\dots\;q_s
\end{equation*}
dan geldt $r=s$ en, na herschikken, is elke $q_i$ geassocieerd met $p_i$:
\begin{equation*}
    q_i = u_i p_i
\end{equation*}
voor een eenheid $u_i$.

\textbf{\underbar{Voorbeeld:}} in $\mathbb{Z}$
\begin{equation*}
    15 = 3 \cdot 5 = (-3) \cdot (-5)
\end{equation*}
Dit zijn dezelfde factorisaties op eenheden na, omdat $-1$ een eenheid is in $\mathbb{Z}$.

\subsection{Resklassenringen \texorpdfstring{$\mathbb{Z}_n$}{Zn} en de eenhedengroep \texorpdfstring{$\mathbb{Z}_n^{\times}$}{Zn*}}

\SimpleHeading{Restklassen modulo $n$}

De ring $\mathbb{Z}_n$ bestaat uit de restklassen modulo $n$:
\begin{equation*}
    \mathbb{Z}_n = \{\bar{0},\bar{1},\dots,\overline{n-1}\}
\end{equation*}
De restklasse van $a$ is
\begin{equation*}
    \bar{a} = \{a +kn \mid k \in \mathbb{Z}\}
\end{equation*}
Twee gehele getallen $a$ en $b$ stellen dezelfde restklasse voor als
\begin{equation*}
    a \equiv b \pmod{n}
\end{equation*}
of equivalent:
\begin{equation*}
    n \mid (a - b)
\end{equation*}
Optelling en vermenigvuldiging zijn gedefinieerd door:
\begin{equation*}
    \begin{aligned}
        \bar{a} + \bar{b}     &= \overline{a+b} \\
        \bar{a} \cdot \bar{b} &= \overline{ab}
    \end{aligned}
\end{equation*}
Deze definities zijn goed gedefinieerd; ze hangen niet af van de gekozen representanten.

\subsubsection{De eenhedengroep \texorpdfstring{$\mathbb{Z}_n^{\times}$}{Znx}}

Niet elk element van $\mathbb{Z}_n$ is inverteerbaar. Bijvoorbeeld in $\mathbb{Z}_{12}$:
\begin{equation*}
    \bar{2} \cdot \bar{6} = \bar{12} = \bar{0}
\end{equation*}
terwijk $\bar{2}$ en $\bar{6}$ niet nul zijn. Ze kunnen dus niet inverteerbaar zijn. \textbf{Een nuldeler kan nooit een eenheid (inverteerbaar element) zijn.}

\begin{nicetheorem}*[Wie is inverteerbaar?]
    \nextpart[]*
    Voor $a \in \mathbb{Z}$ geldt:
    \begin{equation*}
        \bar{a} \in \mathbb{Z}_n^{\times} \iff \operatorname{ggd}(a,\ n) = 1
    \end{equation*}
\end{nicetheorem}

\inlinetext{\underbar{Bewijs:}}[Goldenrod!44] 

\inlinetext{$\Leftarrow$}[Goldenrod!44] Stel eerst dat $\operatorname{ggd}(a,n)=1$. Volgens Bézout bestaan er dan $s,t \in \mathbb{Z}$ zodat:
\begin{equation*}
    1 = sa + tn
\end{equation*} 
Gaan we nu over naar restklassen modulo $n$, dan krijgen we:
\begin{equation*}
    \bar{1} = \overline{sa + tn}
\end{equation*}
Omdat $tn$ deelbaar is door $n$, geldt: $\overline{tn} = \bar{0}$. Dus:
\begin{equation*}
    \bar{1} = \overline{sa} = \bar{s} \cdot \bar{a}
\end{equation*}
Bijgevolg bestaat er dus een element $\bar{s} \in \mathbb{Z}_n$ zodat
\begin{equation*}
    \bar{s} \cdot \bar{a} = \bar{1}
\end{equation*}
Dus $\bar{a}$ is inverteerbaar in $\mathbb{Z}_n$, met inverse $\bar{s}$. Daarom geldt dus: $\bar{a} \in \mathbb{Z}_n^{\times}$.

\inlinetext{$\Rightarrow$}[Goldenrod!44] Stel nu omgekeerd dat $\bar{a} \in \mathbb{Z}_n^{\times}$. Dan bestaat er een $\bar{s} \in \mathbb{Z}_n$ zodat
\begin{equation*}
    \bar{s} \cdot \bar{a} = \bar{1}
\end{equation*}
Dit betekent $\overline{sa} = \bar{1}$. Dus:
\begin{equation*}
    \begin{aligned}
        sa &\equiv 1 \pmod n \\
        &\Downarrow \\
        n &\mid sa - 1 \\
    \end{aligned}
\end{equation*}
Er bestaat dus een $t \in \mathbb{Z}$ zodat $sa - 1 = tn$. Dus:
\begin{equation*}
    1 = sa - tn
\end{equation*}
Neem nu een willekeurige deler $d$ van $a$ en $n$. Dan geldt:
\begin{equation*}
    d \mid a \qquad \text{en} \qquad d \mid n
\end{equation*}
Dan deelt $d$ ook elke lineare combinatie van $a$ en $n$, dus in het bijzonder:
\begin{equation*}
    d \mid (sa - tn)
\end{equation*}
Maar: $sa - tn = 1$, dus
\begin{equation*}
    d \mid 1
\end{equation*}
De enige positieve deler van $1$ is $1$ zelf. Daarom geldt:
\begin{equation*}
    \operatorname{ggd}(a,n) = 1
\end{equation*}

\newpage
\begin{nicedefinition}*[Eenhedengroep]
    \nextpart[]*
    De inverteerbare elementen vormen de \textcolor{PineGreen!75!black}{\textbf{eenhedengroep}}:
    \begin{equation*}
        \mathbb{Z}_n^{\times} = \{\bar{a} \mid \operatorname{ggd}(a,\ n) = 1\}
    \end{equation*}
\end{nicedefinition}

\textbf{\underline{Voorbeeld:}} Zoek $\bar{5}^{-1}$ in $\mathbb{Z}_{12}$.\newline
\underline{Oplossing:} We zoeken eerst $\operatorname{ggd}(12,5)$:
\begin{equation*}
    \begin{aligned}
        12 &= 2 \times 5 + 2 \\
        5  &= 2 \times 2 + 1 \\
        2  &= 2 \times 1 + 0
    \end{aligned}
\end{equation*}
Dus $\operatorname{ggd}(12,5)=1$. De inverse bestaat dus in $\mathbb{Z}_{12}$. Want we hebben net gezien dat een inverse enkel bestaat als $\operatorname{ggd}(5,12)=1$. We doen nu terugsubstitutie:
\begin{equation*}
    \begin{aligned}
        1 &= 1 \times 5 - 2 \times \textcolor{NavyBlue}{2} \\
          &= 1 \times 5 - 2 \times (\textcolor{NavyBlue}{12 - 2 \times 5}) \\
          &= \underbrace{(-2) \times 12}_{\bar{0}} + 5 \times 5
    \end{aligned}
\end{equation*}
Dus $s=5$. Bijgevolg is $\bar{5}^{-1} = \bar{5}$ in $\mathbb{Z}_{12}$.

\subsection{Euler-\texorpdfstring{$\varphi$}{phi} en de stelling van Euler}

\begin{nicedefinition}*[Euler-$\varphi$ functie]
    \nextpart[]*
    Voor $n\in \mathbb{Z}_{\geq 1}$ definiëren we de \textcolor{PineGreen!75!black}{\textbf{Euler-$\varphi$ functie}} door:
    \begin{equation*}
        \varphi(n) = \#\mathbb{Z}_n^{\times}
    \end{equation*}
    Dus $\varphi(n)$ telt het aantal inverteerbare elementen in $\mathbb{Z}_n$. Omdat:
    \begin{equation*}
        \bar{a} \in \mathbb{Z}_n^\times \iff \operatorname{ggd}(a,n) = 1
    \end{equation*}
    geldt ook:
    \begin{equation*}
        \varphi(n) = \#\{a \in \{1, \dots, n\} \mid \operatorname{ggd}(a,n) = 1\}
    \end{equation*}
\end{nicedefinition}

\textbf{\underbar{Intuïtief:}} $\varphi(n)$ telt hoeveel getallen tussen $1$ en $n$ relatief priem zijn. Met \inlinetext{relatief priem} bedoelen we: $a$ is relatief priem met $n$ als:
\begin{equation*}
    \operatorname{ggd}(a,n) = 1
\end{equation*}

\SimpleHeading{Basisformules}

\textbf{\underbar{Geval 1: $p$ priem}}\newline
Als $p$ priem is, dan:
\begin{equationbox}*[NavyBlue][SkyBlue!10]
    \varphi(p) = p-1
\end{equationbox}
Waarom? De getallen $1,2,\dots,p-1$ zijn allemaal priem met $p$, want $p$ heeft geen delers behalve $1$ en $p$. Alleen $p$ is niet zelf relatief priem met $p$, want $\operatorname{ggd}(p,p) = p$.

\textbf{\underbar{Geval 2: macht van een priemgetal}}\newline
Als $p$ priem is en $r \geq 1$, dan:
\begin{equationbox}*[NavyBlue][SkyBlue!10]
    \varphi(p^r) = p^r - p^{r-1}
\end{equationbox}
Ook:
\begin{equationbox}*[NavyBlue][SkyBlue!10]
    \varphi(p^r) = p^{r-1}(p-1)
\end{equationbox}
Waarom? We tellen de getallen tussen $1$ en $p^r$ die relatief priem zijn met $p^r$. Een getal is niet relatief priem met $p^r$ precies wanneer het deelbaar is door $p$.

Dus we nemen alle getallen:
\begin{equation*}
    1,2,\dots,p^r
\end{equation*}
en trekken de veelvouden van $p$ er van af. Er zijn $p^r$ getallen in totaal. De veelvouden van $p$ zijn:
\begin{equation*}
    p, 2p, 3p, \dots, p^{r-1}p = 1
\end{equation*}
Dat zijn er $p^{r-1}$. Dus:
\begin{equation*}
    \varphi(p^r) = p^r - p^{r-1}
\end{equation*}
Factoriseer:
\begin{equation*}
    \varphi(p) = p^{r-1}(p-1)
\end{equation*}

\begin{niceproperty}*[Multiplicativiteit van $\varphi$]
    \nextpart[]*
    \begin{equation*}
        \operatorname{ggd}(m,\ n) = 1 \implies \varphi(mn) = \varphi(m)\varphi(n)
    \end{equation*}
\end{niceproperty}

\SimpleHeading{Stelling van Euler}

De stelling van Euler is een rechtstreeks gevolg van de stelling van Lagrange.

\begin{nicetheorem}*[Stelling van Euler]\label{l6_stelling-euler}
    \nextpart[]*
    Als $\operatorname{ggd}(a,n)=1$, dan geldt:
    \begin{equation*}
        a^{\varphi(n)} \equiv 1 \pmod n
    \end{equation*}
\end{nicetheorem}

\inlinetext{\underbar{Bewijs:}}[Goldenrod!44] Omdat $\operatorname{ggd}(a,n)=1$ weten we dat:
\begin{equation*}
    \bar{a} \in \mathbb{Z}_n^{\times}
\end{equation*}
Ook weet je dat $\mathbb{Z}_n^{\times}$ een groep is onder vermenigvuldiging modulo $n$, en dat de orde van deze groep gelijk is aan:
\begin{equation*}
    \#\mathbb{Z}_n^{\times} = \varphi(n)
\end{equation*}
Uit Lagrange (\Cref{gevolg:lagrange_2_gevolg}) geldt er in een eindige groep $G$:
\begin{equation*}
    g^{\#G} = e \qquad \forall g \in G
\end{equation*}
We passen dit toe, dus:
\begin{equation*}
    G = \mathbb{Z}_n^{\times} \qquad \text{en} \qquad g = \bar{a}
\end{equation*}
Dan krijg je:
\begin{equation*}
    \begin{aligned}
        \bar{a}^{\#\mathbb{Z}_n^{\times}} &= \bar{1} \\
        &\Downarrow \textcolor{Green}{\;\; \#\mathbb{Z}_n^{\times} = \varphi(n)} \\
        \bar{a}^{\varphi(n)} &= \bar{1} \\
        &\Downarrow \\
        a^{\varphi(n)} &\equiv 1 \pmod n
    \end{aligned}
\end{equation*}

\subsubsection{Kleine stelling van Fermat als speciaal geval}

Als $p$ priem is, dan weten we nu al dat $\varphi(p) = p-1$. Dus Euler geeft:
\begin{equation*}
    a^{\varphi(p)} \equiv 1 \pmod p
\end{equation*}
Daarom:
\begin{equation*}
    a^{p-1} \equiv 1 \pmod p
\end{equation*}
op voorwaarde dat $\operatorname{ggd}(a,p) = 1$. Omdat $p$ priem is, betekent dat hetzelfde als:
\begin{equation*}
    p \nmid a
\end{equation*}

\begin{nicecorollary}*[Kleine stelling van Fermat]
    \nextpart[]*
    Als $p$ priem is en $p \nmid a$, dan geldt:
    \begin{equation*}
        a^{p-1} \equiv 1 \pmod p
    \end{equation*}
\end{nicecorollary}

Euler is dus gewoon Lagrange toegepast op de eenhedengroep.

\newpage
\textbf{\underbar{Voorbeeld:}} Bereken $\varphi(12)$.

\underbar{Oplossing:} Factoriseer:
\begin{equation*}
    12 = 4 \cdot 3 = 2^2 \cdot 3
\end{equation*}
Omdat nu $\operatorname{ggd}(4,3)=1$, mag je schrijven:
\begin{equation*}
    \varphi(12) = \varphi(4)\varphi(3)
\end{equation*}
Nu:
\begin{equation*}
    \begin{aligned}
        \varphi(4) &= \varphi(2^2) = 2^2 - 2^{2-1} = 4 - 2 = 2 \\
        \varphi(3) &= 3 - 1 = 2
    \end{aligned}
\end{equation*}
Dus:
\begin{equation*}
    \varphi(12)=2 \cdot 2 = 4
\end{equation*}
Inderdaad: $\mathbb{Z}_{12}^{\times} = \{\bar{1},\bar{5},\bar{7},\bar{11}\}$.

\textbf{\underbar{Voorbeeld:}} Bereken $\varphi(360)$.

\underbar{Oplossing:} Factoriseer:
\begin{equation*}
    \begin{aligned}
        360 &= 36 \cdot 10 \\ 
            &= 9 \cdot 4 \cdot 2 \cdot 5 \\
            &= 3^2 \cdot 2^3 \cdot 5 
    \end{aligned}
\end{equation*}
Dus:
\begin{equation*}
    \varphi(360) = \varphi(2^3)\varphi(3^2)\varphi(5)
\end{equation*}
Bereken elk stuk:
\begin{equation*}
    \begin{aligned}
        \varphi(2^3) &= 2^3 - 2^3 = 8 - 4 = 4 \\
        \varphi(3^2) &= 3^2 - 3^1 = 9 - 3 = 6 \\
        \varphi(5)   &= 5 - 1 = 4
    \end{aligned}
\end{equation*}
Dus $\varphi(360) = 4 \cdot 6 \cdot 4 = 96$.

\textbf{\underbar{Voorbeeld:}} Bereken $7^{1000} \pmod{12}$

\underbar{Oplossing:} $n=12$ en $a=7$, om de stelling van Euler te mogen gebruiken controleer je eerst:
\begin{equation*}
    \operatorname{ggd}(7,12) = 1
\end{equation*}
Dus Euler mag gebruikt worden. We weten ook dat $\varphi(12)=4$. De stelling van Euler zegt dan:
\begin{equation*}
    7^4 \equiv 1 \pmod{12}
\end{equation*}
Nu geldt ook dat:
\begin{equation*}
    1000 = 4 \cdot 250
\end{equation*}
Dus:
\begin{equation*}
    7^{1000} = 7^{4 \cdot 250} = (7^4)^{250}
\end{equation*}
Modulo $12$ is dan:
\begin{equation*}
    (7^4)^{250} \equiv 1^{250} \equiv 1 \pmod{12}
\end{equation*}
Dus: $7^{1000} \equiv 1 \pmod{12}$

\subsection{Productringen, \texorpdfstring{$\mathbb{Z}_{mn} \cong \mathbb{Z}_m \times \mathbb{Z}_n$}{Zmn tgo Zm x Zn} en de Chinese reststelling}

\subsubsection{Productringen}

Als $A$ en $B$ ringen zijn, dan is het Carthesisch product
\begin{equation*}
    A \times B = \{(a,b) \mid a \in A,\ b \in B\}
\end{equation*}
opnieuw een ring met componentsgewijze bewerkingen:
\begin{equation*}
    \begin{gathered}
        (a,b) + (a',b') = (a + a', b + b') \\
        (a,b)(a',b') = (aa', bb')
    \end{gathered}
\end{equation*}
Het nul-element is
\begin{equation*}
    (0_A, 0_B),
\end{equation*}
en, als beide ringen een eenheid hebben, is het eenheidselement
\begin{equation*}
    (1_A, 1_B)
\end{equation*}

\subsubsection{Isomorfisme \texorpdfstring{$\mathbb{Z}_{mn} \cong \mathbb{Z}_m \times \mathbb{Z}_n$}{Zmn tgo Zm x Zn}}

\begin{nicetheorem}*[Splitsingsisomorfisme]
    \nextpart[]*
    Als $m$ en $n$ onderling ondeelbaar zijn, dus
    \begin{equation*}
        \operatorname{ggd}(m,n) = 1
    \end{equation*}
    dan is
    \begin{equation*}
        \psi : \mathbb{Z}_{mn} \to \mathbb{Z}_m \times \mathbb{Z}_n : [a]_{mn} \mapsto ([a]_m,\ [a]_n)
    \end{equation*}
    met 
    \begin{equation*}
        \psi([a]_{mn}) = ([a]_m,\ [a]_n)
    \end{equation*}
    een \textcolor{red!75!black}{\textbf{ringisomorfisme}}.    
\end{nicetheorem}

\textbf{\underbar{Notatie:}} $[a]$ betekent hier gewoon restklasse of equivalentieklasse van $a$ modulo een bepaald getal:
\begin{equation*}
    [a]_n = \bar{a} \in \mathbb{Z}_n = \{a +kn \mid k \in \mathbb{Z}\}
\end{equation*}

\newpage
\inlinetext{\underbar{Bewijs:}}[Goldenrod!44]\newline
\textbf{1) $\psi$ is goed gedefinieerd:} Stel dat $[a]_{mn} = [b]_{mn}$. Dan geldt:
\begin{equation*}
    mn \mid (a-b)
\end{equation*}
Omdat $m \mid mn$ en $n \mid mn$, volgt:
\begin{equation*}
    m \mid (a-b) \qquad \text{en} \qquad n \mid (a-b)
\end{equation*}
Dus $[a]_m = [b]_m$ en $[a]_n = [b]_n$. Bijgevolg is
\begin{equation*}
    ([a]_m, [a]_n)=([b]_m, [b]_n)
\end{equation*}
Dus $\psi$ is goed gedefinieerd.

\textbf{2) $\psi$ is een ringhomomorfisme:} Neem $[a]_{mn},[b]_{mn}\in\mathbb{Z}_{mn}$. Voor de optelling geldt:
\begin{equation*}
    \psi([a]_{mn}+[b]_{mn}) = \psi([a+b]_{mn})
\end{equation*}
Per definitie van $\psi$:
\begin{equation*}
    \psi([a+b]_{mn}) = ([a+b]_m,[a+b]_n)
\end{equation*}
Maar:
\begin{equation*}
    [a+b]_m=[a]_m+[b]_m \qquad \text{en} \qquad [a+b]_n=[a]_n+[b]_n
\end{equation*}
Dus:
\begin{equation*}
    \psi([a]_{mn}+[b]_{mn}) = ([a]_m+[b]_m,[a]_n+[b]_n)
\end{equation*}
Dit is precies:
\begin{equation*}
    ([a]_m,[a]_n)+([b]_m,[b]_n)
\end{equation*}
Dus:
\begin{equation*}
    \psi([a]_{mn}+[b]_{mn}) = \psi([a]_{mn})+\psi([b]_{mn})
\end{equation*}
Voor de vermenigvuldiging geldt analoog:
\begin{equation*}
    \psi([a]_{mn}\cdot[b]_{mn}) = \psi([ab]_{mn}) = ([ab]_m,[ab]_n)
\end{equation*}
Omdat:
\begin{equation*}
    [ab]_m=[a]_m[b]_m \qquad \text{en} \qquad [ab]_n=[a]_n[b]_n
\end{equation*}
krijgen we:
\begin{equation*}
    \psi([a]_{mn}\cdot[b]_{mn}) = ([a]_m[b]_m,[a]_n[b]_n)
\end{equation*}
Dus:
\begin{equation*}
    \psi([a]_{mn}\cdot[b]_{mn}) = \psi([a]_{mn})\cdot\psi([b]_{mn})
\end{equation*}
Ten slotte:
\begin{equation*}
    \psi([1]_{mn})=([1]_m,[1]_n)
\end{equation*}
dus $\psi$ bewaart ook het eenheidselement. Daarom is $\psi$ een ringhomomorfisme.

\textbf{3) $\psi$ is injectief:} We bekijken de kern van $\psi$ als homomorfisme van additieve groepen. Stel:
\begin{equation*}
    [a]_{mn}\in\ker(\psi)
\end{equation*}
Dan:
\begin{equation*}
    \psi([a]_{mn})=([0]_m,[0]_n)
\end{equation*}
Dus:
\begin{equation*}
    ([a]_m,[a]_n)=([0]_m,[0]_n)
\end{equation*}
Hieruit volgt:
\begin{equation*}
    [a]_m=[0]_m \qquad \text{en} \qquad [a]_n=[0]_n
\end{equation*}
Dus:
\begin{equation*}
    m\mid a - \textcolor{Aquamarine}{0} \qquad \text{en} \qquad n\mid a - \textcolor{Aquamarine}{0}
\end{equation*}
Omdat $\operatorname{ggd}(m,n)=1$, volgt: $mn\mid a$. Dus:
\begin{equation*}
    [a]_{mn}=[0]_{mn}
\end{equation*}
Daarom:
\begin{equation*}
    \ker(\psi)=\{[0]_{mn}\}
\end{equation*}
Dus $\psi$ is injectief.

\begin{nicenote}*
    \nextpart[]*
    Om injectiviteit van een groepshomomorfisme te bewijzen, volstaat het om te tonen dat de kern alleen uit het neutraal element bestaat.
\end{nicenote}

\textbf{4) $\psi$ is bijectief:} We hebben: $\#\mathbb{Z}_{mn}=mn$. Ook hebben we:
\begin{equation*}
    \#(\mathbb{Z}_m\times\mathbb{Z}_n) = \#\mathbb{Z}_m\cdot \#\mathbb{Z}_n = mn
\end{equation*}
Dus domein en codomein hebben evenveel elementen. Omdat $\psi$ injectief is tussen twee eindige verzamelingen met hetzelfde aantal elementen, is $\psi$ ook surjectief. Dus $\psi$ is bijectief.

\textbf{Conclusie:} De afbeelding $\psi$ is een bijectief ringhomomorfisme. Daarom is $\psi$ een ringisomorfisme.
Dus: 
\begin{equation*}
    \mathbb{Z}_{mn}\cong \mathbb{Z}_m\times \mathbb{Z}_n
\end{equation*}

\subsubsection{Gevolgen voor \texorpdfstring{$\mathbb{Z}_n^{\times}$}{Znx} en \texorpdfstring{$\varphi$}{phi}}

Uit $\mathbb{Z}_{mn}\cong \mathbb{Z}_m\times \mathbb{Z}_n$ volgt voor de eenhedengroepen:
\begin{equation*}
    \mathbb{Z}_{mn}^\times\cong \mathbb{Z}_m^\times\times \mathbb{Z}_n^\times
\end{equation*}
\textbf{als $\operatorname{ggd}(m,\ n)=1$}. Dus
\begin{equationbox}*[NavyBlue][SkyBlue!10]
    \varphi(mn) = \varphi(m)\varphi(n)
\end{equationbox}
voor onderling ondeelbare $m$ en $n$. Als
\begin{equation*}
    n = p_1^{r_1} p_2^{r_2} \dots p_s^{r_s}
\end{equation*}
met \textbf{verschillende priemgetallen} $p_i$, dan volgt:
\begin{equationbox}*[NavyBlue][SkyBlue!10]
    \varphi(n) = \prod_{i=1}^{s}\varphi(p_i^{r_s})
\end{equationbox}
en dus:
\begin{equationbox}*[NavyBlue][SkyBlue!10]
    \varphi(n) = \prod_{i=1}^{s}p_i^{r_i-1}(p_i -1)
\end{equationbox}
Equivalent:
\begin{equationbox}*[NavyBlue][SkyBlue!10]
    \varphi(n) = n\prod_{p\mid n}^{}\left(1 - \frac{1}{p}\right)
\end{equationbox}

\subsubsection{Chinese reststelling}

De Chinese reststelling hangt rechtstreeks samen met het ringisomorfisme
\begin{equation*}
    \mathbb{Z}_{mn}\cong \mathbb{Z}_m\times \mathbb{Z}_n
\end{equation*}
wanneer
\begin{equation*}
    \operatorname{ggd}(m,n)=1.
\end{equation*}
De surjectiviteit van dit isomorfisme betekent precies dat elke combinatie van een rest modulo $m$ en een rest modulo $n$ afkomstig is van één rest modulo $mn$.

\begin{nicetheorem}*[Chinese reststelling: twee moduli]
    \nextpart[]*
    Zij $m,n\in\mathbb{Z}_{\geq 1}$ met
    \begin{equation*}
        \operatorname{ggd}(m,n)=1.
    \end{equation*}

    Dan bestaat er voor elke keuze van $a_1,a_2\in\mathbb{Z}$ een $a\in\mathbb{Z}$
    zodat
    \begin{equation*}
        a\equiv a_1 \pmod m
    \end{equation*}
    en
    \begin{equation*}
        a\equiv a_2 \pmod n.
    \end{equation*}

    Bovendien is deze oplossing uniek modulo $mn$.
\end{nicetheorem}

\SimpleHeading{Betekenis}

We kunnen dus altijd een systeem van twee congruenties oplossen:
\begin{equation*}
    \begin{cases}
        a\equiv a_1 \pmod m,\\
        a\equiv a_2 \pmod n,
    \end{cases}
\end{equation*}
zolang $m$ en $n$ onderling ondeelbaar zijn. De oplossing is uniek modulo $mn$. Dat betekent: als $a$ en $a'$ allebei oplossingen zijn, dan geldt
\begin{equation*}
    a\equiv a' \pmod{mn}
\end{equation*}

\SimpleHeading{Expliciete formule via Bézout}

Omdat $\operatorname{ggd}(m,n)=1$, bestaan er volgens Bézout gehele getallen $s,t\in\mathbb{Z}$ zodat
\begin{equation*}
    sm + tn = 1
\end{equation*}
Dan is een oplossing van
\begin{equation*}
    \begin{cases}
        a\equiv a_1 \pmod m,\\
        a\equiv a_2 \pmod n
    \end{cases}
\end{equation*}
gegeven door:
\begin{equationbox}*[NavyBlue][SkyBlue!10]
    a=sma_2+tna_1
\end{equationbox}

\textbf{\underline{Controle modulo $m$:}}

Modulo $m$ geldt: $sm\equiv 0 \pmod m$. Dus $sma_2\equiv 0 \pmod m$. Uit $sm+tn=1$ volgt modulo $m$:
\begin{equation*}
    tn\equiv 1 \pmod m
\end{equation*}
Daarom:
\begin{equation*}
    tna_1\equiv a_1 \pmod m
\end{equation*}
Dus:
\begin{equation*}
    a=sma_2+tna_1\equiv a_1 \pmod m
\end{equation*}

\textbf{\underline{Controle modulo $n$:}}

Modulo $n$ geldt: $tn\equiv 0 \pmod n$. Dus $tna_1\equiv 0 \pmod n$. Uit $sm+tn=1$ volgt modulo $n$:
\begin{equation*}
    sm\equiv 1 \pmod n
\end{equation*}
Daarom:
\begin{equation*}
    sma_2\equiv a_2 \pmod n
\end{equation*}
Dus:
\begin{equation*}
    a=sma_2+tna_1\equiv a_2 \pmod n
\end{equation*}
Hiermee is gecontroleerd dat $a=sma_2+tna_1$ inderdaad een oplossing is.

\SimpleHeading{Oplossing is uniek}

Stel dat $a$ en $a'$ allebei oplossingen zijn. Dan geldt:
\begin{equation*}
    a\equiv a' \pmod m \qquad \text{en} \qquad a\equiv a' \pmod n
\end{equation*}
Dus:
\begin{equation*}
    m\mid(a-a') \qquad \text{en} \qquad n\mid(a-a')
\end{equation*}
Omdat $\operatorname{ggd}(m,n)=1$, volgt dat:
\begin{equation*}
    mn\mid(a-a')
\end{equation*}
Dus:
\begin{equation*}
    a\equiv a' \pmod{mn}
\end{equation*}
De oplossing is dus uniek modulo $mn$

\SimpleHeading{Meerder moduli}

De Chinese reststelling bestaat ook voor meer dan twee moduli. Zij
\begin{equation*}
    m_1,m_2,\dots,m_r
\end{equation*}
paarsgewijs onderling ondeelbaar. Dat betekent:
\begin{equation*}
    \operatorname{ggd}(m_i,m_j)=1 \qquad \text{voor alle } i\neq j
\end{equation*}
Dan heeft het systeem
\begin{equation*}
    \begin{cases}
        x\equiv a_1 \pmod {m_1},\\
        x\equiv a_2 \pmod {m_2},\\
        \vdots\\
        x\equiv a_r \pmod {m_r}
    \end{cases}
\end{equation*}
een oplossing. Bovendien is die oplossing uniek modulo
\begin{equation*}
    M=m_1m_2\cdots m_r
\end{equation*}
Met andere woorden: alle oplossingen verschillen van elkaar met een veelvoud van $M$.

\textbf{\underbar{Voorbeeld:}} Los het volgende systeem op:
\begin{equation*}
    \begin{cases}
            x\equiv 2 \pmod 3,\\
            x\equiv 3 \pmod 5.
        \end{cases}
\end{equation*}
Hier is
\begin{equation*}
    m=3,\qquad n=5,\qquad a_1=2,\qquad a_2=3
\end{equation*}
Omdat $\operatorname{ggd}(3,5)=1$, mogen we de Chinese reststelling toepassen. We zoeken eerst $s,t\in\mathbb{Z}$ zodat
\begin{equation*}
    sm+tn=1
\end{equation*}
Dus:
\begin{equation*}
    3s+5t=1
\end{equation*}
Een oplossing is:
\begin{equation*}
    2\cdot 3+(-1)\cdot 5=1
\end{equation*}
Dus
\begin{equation*}
    s=2 \qquad \text{en} \qquad t=-1
\end{equation*}
De expliciete formule uit de Chinese reststelling is:
\begin{equation*}
    x=sma_2+tna_1
\end{equation*}
Invullen geeft:
\begin{equation*}
    x=(2\cdot 3)\cdot 3+((-1)\cdot 5)\cdot 2
\end{equation*}
Dus $x=18-10=8$. We controleren: $8\equiv 2 \pmod 3$, want $8=2\cdot 3+2$.

Ook: $8\equiv 3 \pmod 5$, want $8=1\cdot 5+3$. Dus $x=8$ is een oplossing. Omdat de oplossing uniek is modulo
\begin{equation*}
    mn=3\cdot 5=15
\end{equation*}
is de volledige oplossing:
\begin{equation*}
    x\equiv 8 \pmod {15}
\end{equation*}

\subsection{Eindige velden}

Als $p$ priem is, dan is
\begin{equation*}
    \mathbb{F}_p = \mathbb{Z}_p
\end{equation*}
een veld met $p$ elementen.

\textbf{\underbar{Waarom?}} Neem een niet-nul element $\bar{a} \in \mathbb{Z}_p$. Omdat $p$ priem is en $a \not\equiv 0 \pmod p$, geldt:
\begin{equation*}
    \operatorname{ggd}(a,p) = 1
\end{equation*}
Door Bézout bestaan er dan $s,\ t$ zodat
\begin{equation*}
    sa + tp = 1
\end{equation*}
Modulo $p$ wordt dit:
\begin{equation*}
    \bar{s}\bar{a} = \bar{1}
\end{equation*}
Dus $\bar{a}$ heeft een inverse. Elk niet-nul element is dus inverteerbaar. Dus $\mathbb{Z}_p$ is een veld.

\subsubsection{Velden construeren met veeltermen modulo een irreduceerbare veelterm}

\SimpleHeading{Setup}

Laat $F$ een veld zijn en beschouw de Euclidische ring $F[x]$. Neem dan een irreduceerbare veelterm
\begin{equation*}
    g(x) \in F[x]
\end{equation*}
We beschouwen de quotiëntring
\begin{equation*}
    F[x]/(g(x))
\end{equation*}
Dit wordt in de literatuur ook wel lichter genoteerd als 
\begin{equation*}
    F[x] \;\text{ mod } g(x)
\end{equation*}
Hier is $(g(x))$ het ideaal van alle veelvouden van $g(x)$:
\begin{equation*}
    (g(x)) = \{h(x)g(x) \mid h(x) \in F[x]\}
\end{equation*}

\SimpleHeading{Equivalentie modulo $g(x)$}

Twee veeltermen $f_1(x)$ en $f_2(x)$ stellen dezelfde klasse voor modulo $g(x)$ als
\begin{equation*}
    f_1(x) \equiv f_2(x) \pmod{g(x)}
\end{equation*}
Dit betekent:
\begin{equation*}
    g(x) \mid (f_1(x) - f_2(x))
\end{equation*}
Dit is volledig analoog aan rekenen modulo $n$ in $\mathbb{Z}$.

\subsubsection{Waarom \texorpdfstring{$F[x]/(g(x))$}{F[x]/(g(x))} een veld is als \texorpdfstring{$g$}{g} irreduceerbaar is}

\begin{nicetheorem}*
    \nextpart[]*
    Als $g$ irreduceerbaar is, dan is $F[x]/(g(x))$ een \textbf{veld}.
\end{nicetheorem}

\inlinetext{\underline{Bewijs:}}[Goldenrod!44] We moeten aantonen dat elk niet-nul element een inverse heeft. Neem een niet-nul klasse
\begin{equation*}
    [f(x)] \in F[x]/(g(x))
\end{equation*}
Niet-nul betekent:
\begin{equation*}
    g(x) \nmid f(x)
\end{equation*}
want als $f(x)$ een veelvoud is van $g(x)$ is het de klasse $\bar{0} = [f(x)]$. Omdat $g(x)$ irreduceerbaar is (aanname), zijn de enige delers van $g(x)$ eenheden en geassocieerden van $g(x)$. Dus als $g(x)$ geen deler is van $f(x)$, dan moet wel gelden:
\begin{equation*}
    d(x) = \operatorname{ggd}(f(x),\ g(x)) = 1
\end{equation*}
Want als de ggd een geassocieerde van $g(x)$ zou zijn, dan bestaat er een:
\begin{equation*}
    d(x) = c \cdot g(x) \;\;\xLongrightarrow{g\text{ niet-nul}}\;\; g(x) = \frac{1}{c} \cdot d(x)
\end{equation*}
Maar $d(x) \mid f(x)$, dus er bestaat ook een $k(x)$ zodat:
\begin{equation*}
    f(x) = k(x) \cdot d(x) \;\;\xLongrightarrow{d(x)\text{ invullen}}\;\; \begin{aligned}
        f(x) &= k(x) \cdot (c \cdot g(x)) \\
             &= (c \cdot k(x)) \cdot \textcolor{Aquamarine}{g(x)}
    \end{aligned}
\end{equation*} 
Dus hier staat dat $g(x) \mid f(x)$, maar dan kan niet volgens onze aanname!

Door Bézout in de Euclidische ring $F[x]$ bestaan veeltermen $s(x), t(x)$ zodat
\begin{equation*}
    s(x)g(x) + t(x)f(x) = 1
\end{equation*}
Modulo $g(x)$ verdwijnt de eerste term:
\begin{equation*}
    [t(x)][f(x)] = [1]
\end{equation*}
Dus $[t(x)]$ is de inverse van $[f(x)]$. Daarom is 
\begin{equation*}
    F[x]/(g(x))
\end{equation*}
een veld!

\begin{barbox}[WildStrawberry][gray!3]
    \textcolor{WildStrawberry}{\textbf{Het parallel:}} een \textbf{priemgetal} $p$ maakt $\mathbb{Z}/(p)$ een veld; een \textbf{irreduceerbare veelterm} $g$ maakt $F[x]/(g)$ een veld. Irreduceerbare veeltermen zijn de "priemgetallen" van $F[x]$. 
\end{barbox}

\textbf{\underbar{Voorbeeld:}} We construeren het eindige veld met $4$ elementen (genoteerd als $\mathbb{F}_{2^2}$ of $\mathbb{F}_4$) door de veeltermring $\mathbb{F}_2[x]$ te delen door een irreduceerbare veelterm van graad $2$.

\textbf{1) Keuze irreduceerbare veelterm $g(x)$:} We kiezen de veelterm
\begin{equation*}
    g(x) = x^2 + x + 1 \in \mathbb{F}_2[x]
\end{equation*}
en checken op irreduceerbaarheid. Omdat $g(x)$ graad $2$ heeft, is hij irreduceerbaar \emph{asa} hij geen nulpunten heeft in $\mathbb{F}_2 = \{0,1\}$:
\begin{equation*}
    \begin{aligned}
        g(0) &= 0^2 + 0 + 1 = 1 \neq 0 \\
        g(1) &= 1^2 + 1 + 1 = 3 \equiv 1 \pmod 2 \neq 0
    \end{aligned}
\end{equation*}
Omdat er geen nulpunten zijn, is $g(x)$ \textbf{irreduceerbaar over $\mathbb{F}_2$}. De quotiëntring $\mathbb{F}[x]/(x^2 + x + 1)$ is hierdoor gegarandeerd een \textbf{veld}.

\textbf{2) De elementen van het veld:} Als we een willekeurige veelterm delen door $g(x)$ (graad 2), zijn de mogelijke resten $r(x)$ altijd van graad $< 2$. De algemene vorm is: 
\begin{equation*}
    r(x) = ax + b \pmod{g(x)}
\end{equation*}
waarbij de coëfficiënten $a$ en $b$ uit het basisveld $\mathbb{F}_2 = \{0,1\}$ moeten komen. Dit levert exact $2^2=4$ unieke combinaties (elementen) op:
\begin{itemize}
    \item $a=0,\ b=0 \to 0$
    \item $a=0,\ b=1 \to 1$
    \item $a=1,\ b=0 \to x$
    \item $a=1,\ b=1 \to x+1$
\end{itemize}
Het veld is dus de verzameling:
\begin{equation*}
    \mathbb{F}_2[x]/(x^2+x+1) = \{0,\ 1,\ x,\ x+1\}
\end{equation*}

\textbf{3) Rekenen in dit veld:} Bij het rekenen in dit veld gelden er twee belangrijke regels:
\begin{enumerate}
    \item \textbf{Coëfficiënten:} Rekenen met getallen gaat modulo $2$
    \item \textbf{Machten van $x$:} Rekenen met $x$ gaat modulo $x^2 + x + 1$. Dus $x^2 + x + 1 = \bar{0}$
\end{enumerate}

\begin{equation*}
\begin{array}{c@{\hspace{2cm}}c}
    \begin{array}{c|cccc}
        + & 0 & 1 & x & x+1 \\
        \hline
        0     & 0   & 1   & x   & x+1 \\
        1     & 1   & 0   & x+1 & x \\
        x     & x   & x+1 & 0   & 1 \\
        x+1   & x+1 & x   & 1   & 0
    \end{array}
    &
    \begin{array}{c|cccc}
        \cdot & 0 & 1 & x & x+1 \\
        \hline
        0     & 0 & 0   & 0   & 0 \\
        1     & 0 & 1   & x   & x+1 \\
        x     & 0 & x   & x+1 & 1 \\
        x+1   & 0 & x+1 & 1   & x
    \end{array}
\end{array}
\end{equation*}

\begin{itemize}
    \item $x^2 + x + 1 = 0 \implies x^2 = x + 1$
    \item $x \cdot x = x^2 = x + 1$
    \item $x \cdot (x + 1) = x^2 + x = x + 1 + x = 2x + 1 = 1$
\end{itemize}

Merk ook op dat elk niet-nul element een inverse heeft, het is dus zeker een veld!

\subsubsection{Algemene eindige velden \texorpdfstring{$\mathbb{F}_{p^m}$}{Fpm}}

We combineren nu twee ideeën:
\begin{enumerate}
    \item $\mathbb{F}_p$ is een veld met $p$ elementen.
    \item Als $g(x)$ irreduceerbaar is over een veld, dan is $F[x]/\left(g(x)\right)$ opnieuw een veld.
\end{enumerate}
Neem dus een irreduceerbare veelterm
\begin{equation*}
    g(x) \in \mathbb{F}_p[x]
\end{equation*}
met
\begin{equation*}
    \text{deg}(g) = m
\end{equation*}
Dan is
\begin{equationbox}*[NavyBlue][SkyBlue!10]
    \mathbb{F}_p[x]/\left(g(x)\right)
\end{equationbox}
een \textbf{veld!}

\SimpleHeading{Aantal elementen}

Door Euclidische deling heeft elke klasse modulo $g(x)$ een unieke representant met graad kleiner dan $m$:
\begin{equation*}
    a_0 + a_1x + a_2x^2 + \cdots + a_{m-1}x^{m-1}
\end{equation*}
waarbij
\begin{equation*}
    a_i \in \mathbb{F}_p
\end{equation*}
Er zijn $m$ coëfficiënten, en elke coëfficiënt heeft $p$ mogelijke waarden. Daarom zijn er
\begin{equationbox}*[NavyBlue][SkyBlue!10]
    p^m
\end{equationbox}
elementen. Dus:
\begin{equationbox}*[NavyBlue][SkyBlue!10]
    \#\left(\mathbb{F}_p[x]/\left(g(x)\right)\right) = p^m
\end{equationbox}
Dit veld noteert men als
\begin{equation*}
    \mathbb{F}_{p^m} \qquad \text{of} \qquad GF(p^m)
\end{equation*}

\SimpleHeading{Interpretatie}

Rekenen in $GF(p^m)$ betekent:
\begin{enumerate}
    \item Reken met veeltermen met coëfficiënten in $\mathbb{F}_p$.
    \item Reduceer telkens modulo de irreduceerbare veelterm $g(x)$.
    \item Het resultaat kan steeds als veelterm van graad kleiner dan $m$ geschreven worden.
\end{enumerate}

\textbf{\underbar{Voorbeeld:}}\newline
Als $m=3$, dan ziet elk element er uit als
\begin{equation*}
    a_0 + a_1x + a_2x^2 \qquad \text{met} \qquad a_i \in \mathbb{F}_p
\end{equation*}
Als een product een term $x^3$ of hoger oplevert, reduceer je met de relatie $g(x) = 0$.

\subsubsection{Belangrijke feiten over eindige velden}

\begin{enumerate}
    \item Elk eindig veld heeft een aantal elementen van de vorm
    \begin{equation*}
        p^m
    \end{equation*}
    waar $p$ een \textbf{priemgetal} is en $m \ge 1$.
    \item Voor elk priemgetal $p$ en elke $m \ge 1$ bestaat er een irreduceerbare veelterm van graad $m$ over $\mathbb{F}_p$.
    \item Dus voor elke $p$ en $m$ bestaat er een eindig veld met $p^m$ elementen.
    \item Alle velden met $p^m$ elementen zijn \textbf{isomorf}. Daarom mag men spreken over \textbf{het} eindige veld met $p^m$ elementen:
    \begin{equation*}
        \mathbb{F}_{p^m} = GF(p^m)
    \end{equation*}
    In de volgende les wordt er intensief gewerkt met $\mathbb{F}_{2^8}$, wat belangrijk is in toepassingen zoals codetheorie en cryptografie.
\end{enumerate}

\subsection{Belangrijkste algoritmische patronen}

\subsubsection{Inverse modulo \texorpdfstring{$n$}{n}}

Om $\overline{a}^{-1}$ in $\mathbb{Z}_n$ te vinden:
\begin{enumerate}
    \item Controleer dat $\operatorname{ggd}(a,n) = 1$.
    \item Gebruik extended Euclid:
    \begin{equation*}
        sa + tn = 1
    \end{equation*}
    \item Dan:
    \begin{equation*}
        \overline{a}^{-1} = \overline{s}
    \end{equation*}
\end{enumerate}

\subsubsection{Inverse modulo een veelterm}

Om $\left[f(x)\right]^{-1}$ in F$\left[x\right]/\left(g(x)\right)$ te vinden:
\begin{enumerate}
    \item Contoleer dat $g(x)$ irreduceerbaar is en $\left[f(x)\right] \neq [0]$.
    \item Gebruik extended Euclid in $F[x]$:
    \begin{equation*}
        s(x)g(x) = t(x)f(x) = 1
    \end{equation*}
    \item Dan:
    \begin{equation*}
        \left[f(x)\right]^{-1} = \left[t(x)\right]
    \end{equation*}
\end{enumerate}

\subsubsection{Chinese reststelling oplossen}

Om
\begin{equation*}
    x \equiv a_1 \pmod m, \qquad x \equiv a_2 \pmod n
\end{equation*}
met $\operatorname{ggd}(m, n) = 1$ op te lossen:
\begin{enumerate}
    \item Bereken $s,t$ met
    \begin{equation*}
        sm + tn = 1
    \end{equation*}
    \item Neem
    \begin{equation*}
        x = sm\ \textcolor{Aquamarine}{a_2} + tn\ \textcolor{Aquamarine}{a_1}
    \end{equation*}
    \item De oplossing is uniek modulo $mn$.
\end{enumerate}

\subsection{Carmichael}

Deze stelling hebben we in de les niet gezien, maar kan handig zijn.

\begin{nicetheorem}*[Carmichael]\label{l6_carmichael}
    \nextpart[]*
    Voor elk geheel getal $n \ge 1$ bestaat er een kleinste positief geheel getal $\lambda(n)$ zodat voor elk geheel getal $a$ met
    \begin{equation*}
        \operatorname{ggd}(a,\ n) = 1
    \end{equation*}
    geldt:
    \begin{equation*}
        a^{\lambda(n)} \equiv 1 \pmod n
    \end{equation*}
    Het getal $\lambda(n)$ heet de \textcolor{red!75!black}{\textbf{Carmichael-functie}} van $n$.
\end{nicetheorem}

Net zoals de stelling van Euler zijn er rekenregels:

\SimpleHeading{Rekenregel 1 -- Voor een priemgetal $p$}

Uit de Kleine Stelling van Fermat weten we dat als $p$ een priemgetal is en $p \nmid a$, dan geldt:
\begin{equation*}
    a^{p-1} \equiv 1 \pmod p
\end{equation*}
Dit is de stelling van Carmichael als $\lambda(n) = p-1$. Dus er geldt steeds voor een priemgetal dat:
\begin{equationbox}*[NavyBlue][SkyBlue!10]
    \lambda(p) = p - 1
\end{equationbox}

\SimpleHeading{Rekenregel 2 -- product twee onderling ondeelbare getallen}

Als je twee getallen $x$ en $y$ hebt die onderling ondeelbaar zijn ($\operatorname{ggd}(x, y) = 1$), dan geldt de volgende universele rekenregel voor de Carmichael-functie:
\begin{equationbox}*[NavyBlue][SkyBlue!10]
    \lambda(xy) = \operatorname{kgv}\left(\lambda(x),\ \lambda(y)\right)
\end{equationbox}

\inlinetext{\underbar{Bewijs:}}[Goldenrod!66] Stel voor het gemak dat:
\begin{equation*}
    L = \operatorname{kgv}\left(\lambda(x),\ \lambda(y)\right)
\end{equation*}
We bewijzen beide deelbaarheden.

\textbf{1. Bewijs dat $\lambda(xy) \mid L$}\newline
Neem een willekeurig geheel getal $a$ waarvoor geldt dat
\begin{equation*}
    \operatorname{ggd}(a,\ xy) = 1
\end{equation*}
Dan is $a$ zowel onderling ondeelbaar met $x$ als met $y$:
\begin{equation*}
    \operatorname{ggd}(a,x) = 1 \qquad \text{en} \qquad \operatorname{ggd}(a,y) = 1
\end{equation*}
Uit de definitie van Carmichael volgt dan dat:
\begin{equation*}
    a^{\lambda(x)} \equiv 1 \pmod x \qquad \text{en} \qquad a^{\lambda(y)} \equiv \pmod y
\end{equation*}
Ook weet je dat, omdat $\lambda(x) \mid L$, er een geheel getal $r$ bestaat zodat
\begin{equation*}
    L = r \cdot \lambda(x)
\end{equation*}
en dus:
\begin{equation*}
    a^L = \left(a^{\lambda(x)}\right)^r \equiv 1^r \equiv 1 \pmod x
\end{equation*}
Analoog, omdat $\lambda(y) \mid L$, geldt: $a^L \equiv 1 \pmod y$. We hebben dus
\begin{equation*}
    a^L \equiv 1 \pmod x \qquad \text{en} \qquad a^L \equiv 1 \pmod y
\end{equation*}
Omdat $x$ en $y$ onderling ondeelbaar zijn, volgt uit de Chinese reststelling:
\begin{equation*}
    a^L \equiv 1 \pmod {xy}
\end{equation*}
Dus $L$ is een exponent die werkt voor alle getallen die onderling ondeelbaar zijn met $xy$. Omdat $\lambda(xy)$ de \textbf{kleinste} positieve dergelijke exponent is, geldt:
\begin{equationbox}*[black]
    \lambda(xy) \mid L
\end{equationbox}
Alle exponenten die voor alle eenheden modulo $xy$ werken, zijn namelijk veelvouden van $\lambda(xy)$.

\textbf{2. Bewijs dat $L \mid \lambda(xy)$}\newline
Stel $t = \lambda(xy)$. We tonen eerst aan dat
\begin{equation*}
    \lambda(x) \mid t
\end{equation*}
Neem een willekeurige $a$ met $\operatorname{ggd}(a,x) = 1$. We kunnen niet zomaar zeggen dat $a$ ook onderling ondeelbaar is met $y$. Daarom construeren we met de Chinese reststelling een getal $b$ waarvoor
\begin{equation*}
    \begin{cases}
        b \equiv a \pmod x \\
        b \equiv 1 \pmod y
    \end{cases}
\end{equation*}
Zo een $b$ bestaat want $\operatorname{ggd}(x,y) = 1$. Waarom precies deze twee voorwaarden?
\begin{itemize}
    \item Modulo $x$ gedraagt $b$ zich hetzelfde als $a$.
    \item Modulo $y$ is $b$ gelijk aan $1$, dus zeker onderling ondeelbaar met $y$.
\end{itemize}
Meer nog, omdat $b \equiv a \pmod x$ ($\sim$ $a$ en $b$ hebben dezelfde restklasse bij deling door $x$) geldt:
\begin{equation*}
    \begin{aligned}
        b &= a + k \cdot x \\
        &\Updownarrow \\
        a &= b - k \cdot x
    \end{aligned}
\end{equation*}
Maar er geldt ook dat $\operatorname{ggd}(a,x) = 1$, dus:
\begin{equation*}
    \operatorname{ggd}(b, x) = 1
\end{equation*}
want stel dat $b$ en $x$ wél een gemeenschappelijke deler zouden hebben, bijvoorbeeld $g$. Dan is $g$ een deler van $b$ én $x$, en dus:
\begin{itemize}
    \item $b$ is een veelvoud van $g$
    \item $x$ is een veelvoud van $g$
\end{itemize}
Maar kijk naar de bovenstaande formule voor $a$: $a = \textcolor{Aquamarine}{\mathbf{b}} - k \cdot \textcolor{Aquamarine}{\mathbf{x}}$ \newline
$\to$ Als $g$ een deler is van $b$ en $g$, dan moet ze ook wel een deler zijn van $a$, maar we hebben aangenomen dat dat niet kan!

Omdat $b \equiv 1 \pmod y$, geldt ook dat $\operatorname{ggd}(b, y) = 1$. Dus
\begin{equation*}
    \begin{cases}
        \operatorname{ggd}(b, x) &= 1 \\
        \operatorname{ggd}(b, y) &= 1
    \end{cases}\qquad \Longrightarrow \qquad \operatorname{ggd}(b,\ xy) = 1
\end{equation*}
Volgens de definitie van $t = \lambda(xy)$ en de stelling van Carmichael geldt dan
\begin{equation*}
    b^t \equiv 1 \pmod {xy}
\end{equation*}
In het bijzonder:
\begin{equation*}
    b^t \equiv 1 \pmod x
\end{equation*}
want:
\begin{equation*}
    \begin{aligned}
        b^t &\equiv 1 \pmod {xy} \\
        &\Downarrow \\
        b^t &= 1 + c \cdot (xy) \\
        &\Updownarrow \\
        b^t - 1 &= c \cdot (xy) \\
        &\Updownarrow \\
        b^t - 1 &= (cy) \cdot x
    \end{aligned}
\end{equation*}
Dus $(b^t-1)$ is automatisch ook een veelvoud van $x$!

Maar we hebben $b$ zo geconstrueerd dat $b$ zich modulo $x$ hetzelfde gedraagt als $a$ ($b \equiv a \pmod x$). Dus
\begin{equation*}
    a^t \equiv 1 \pmod x
\end{equation*}
Dit geldt voor iedere $a$ die onderling ondeelbaar is met $x$. Bijgevolg is $t$ een universele exponent modulo $x$, zodat
\begin{equation*}
    \lambda(x) \mid t
\end{equation*}
Op dezelfde manier construeren we voor ieder getal dat onderling ondeelbaar is met $y$ een getal dat modulo $x$ gelijk is aan $1$. Daaruit volgt dan:
\begin{equation*}
    \lambda(y) \mid t
\end{equation*}
Omdat dus zowel $\lambda(x)$ als $\lambda(y)$ delers zijn van $t$, is hun kleinste gemeenschappelijke veelvoud ook een deler van $t$:
\begin{equationbox}*
    \operatorname{kgv}\left(\lambda(x), \lambda(y)\right) \mid t 
\end{equationbox}
Dus
\begin{equationbox}*
    L \mid \lambda(xy)
\end{equationbox}

\textbf{3. Conclusie} \newline
We hebben beide deelbaarheden:
\begin{equation*}
    \lambda(xy) \mid L \qquad \text{en} \qquad L \mid \lambda(xy)
\end{equation*}
Daarom zijn beide positieve getallen gelijk:
\begin{equationbox}*
    \lambda(xy) = \operatorname{kgv}\left(\lambda(x), \lambda(y)\right)
\end{equationbox}

\SimpleHeading{Rekenregel 3}

Voor elk positief geheel getal $N$ geldt:

\begin{equationbox}*[NavyBlue][SkyBlue!10]
    \lambda(N) \mid \varphi(N)
\end{equationbox}

\inlinetext{\underbar{Bewijs:}}[Goldenrod!66] Beschouw de eenhedengroep (multiplicatieve groep van de inverteerbare restklassen modulo $N$):
\begin{equation*}
    G = \mathbb{Z}_N^\times = \left(\mathbb{Z}/N\mathbb{Z}\right)^\times
\end{equation*}
Per definitie bevat deze groep precies $\varphi(N)$ elementen:
\begin{equation*}
    \#G = \varphi(N)
\end{equation*}
Volgens Lagrange geldt dan: voor elk element $a \in G$ geldt dat de orde van $a$ de orde van de groep deelt (\Cref{gevolg:H3_gevolg_lagrange_1}):
\begin{equation*}
    \operatorname{ord}(a) \mid \varphi(N)
\end{equation*}
De Carmichael-functie $\lambda(N)$ is de exponent van de groep $G$. Dit betekent dat
\begin{equation*}
    \operatorname{kgv}\left\{\operatorname{ord}(a) \mid a \in G\right\}
\end{equation*}
Elke orde ($\operatorname{ord}(a)$) is dus een deler van $\varphi(N)$. Het kleinste gemene veelvoud van getallen die allemaam $\varphi(N)$ delen, deelt dus zelf ook $\varphi(N)$. Daarom geldt:
\begin{equationbox}
    \lambda(N) \mid \varphi(N)
\end{equationbox}

\subsection{Samenvattende Oefeningen}

\emoji{pushpin} \inlinecode{TODO}[red!22] \emoji{pushpin}










